\section{从复杂度基准到 Agent 评测范式}

上一节提出 population training 的总体命题，本节先说明它为何在任务复杂度上必要。比较 Atari、Go 与 StarCraft 时，应同时看信息可见性、参与者数量、动作空间和时间跨度；这些维度共同决定固定策略和固定评测对手能否代表部署环境。

讲者先通过经典任务对比，解释为什么 StarCraft II 一直被视作现代 Agent 研究的``中间地带''：它既不像 Atari 那样简化，也不像开放互联网任务那样无边界，却已经拥有多主体、部分可观测、长时信用分配和超大动作空间等要素。

\begin{figure}[H]
\centering
\includegraphics[width=0.92\textwidth]{frame-01.jpg}
\caption{复杂度基准对比：Atari、Go 与 StarCraft}
\label{fig:complexity-benchmark}
\end{figure}

\begin{knowledgebox}{读图：复杂度不是一条单轴曲线}
Atari 主要增加感知和控制难度，Go 增加深搜索与对手响应，StarCraft 同时加入不完美信息、多单位控制和超长策略链。LLM Agent 又把规则与用户目标开放化。图支持的是训练分布必须随交互结构扩展，而不是宣称 StarCraft 的动作数可直接等价于自然语言复杂度。
\end{knowledgebox}
图 \ref{fig:complexity-benchmark} 是讲者在开场阶段给出的关键图\footnote{来源：\href{https://www.youtube.com/watch?v=ntjOxjZMaac}{Lecture 03 视频帧}，时间戳 00:05:45。}。从中可以直接读出一个对后文非常重要的结论：当信息不完全、玩家数量上升、动作空间爆炸时，``单策略+固定对手''训练会迅速失效。

\begin{table}[H]
\centering
\caption{从封闭博弈到开放任务的能力需求变化}
\begin{tabular}{p{2.3cm}p{2.9cm}p{2.9cm}p{5.4cm}}
\toprule
任务域 & 状态可观测性 & 对手/用户形态 & 训练关注点 \\
\midrule
Atari & 近似完全可观测 & 单环境反馈 & 样本效率、稳定收敛 \\
Go & 完全可观测 & 双人对弈 & 搜索与策略价值协同 \\
StarCraft II & 部分可观测 & 多策略对手生态 & population diversity、counter-strategy \\
LLM Agent 真实场景 & 开放非结构化 & 协作用户 + 对抗用户 & alignment robustness、tool reliability、攻击面收敛 \\
\bottomrule
\end{tabular}
\end{table}

\begin{importantbox}{评测范式变化}
当任务从``封闭规则''走向``开放交互''，评测指标必须从单一平均分转向多维指标：成功率、最坏分位、恢复能力（recovery）、对抗暴露面（attack surface）、策略多样性与长期稳定性。
\end{importantbox}

\begin{knowledgebox}{为什么讲者强调``不是零和''}
在 LLM 场景里，用户经常是协作者，不一定是对手。这意味着我们需要同时评估 cooperative efficiency 与 adversarial robustness。仅仅复刻 self-play Elo 曲线并不足够。
\end{knowledgebox}

\subsection{本章小结}
本章给出的核心启示是：LLM Agent 的评测必须继承多智能体博弈中的``分布视角''，而不是沿用单任务准确率思路。没有分布鲁棒性，就没有可部署性。

